\documentclass[11pt]{article}
\usepackage[a4paper,margin=0.9in]{geometry}
\usepackage[T1]{fontenc}
\usepackage{helvet}
\renewcommand{\familydefault}{\sfdefault}
\usepackage{xcolor}
\usepackage{booktabs}
\usepackage{array}
\usepackage{colortbl}
\usepackage{float}
\usepackage{amsmath}
\usepackage[font=footnotesize,labelformat=empty]{caption}
\usepackage{titlesec}
\usepackage{fancyhdr}
\usepackage{enumitem}
\usepackage{lastpage}
\usepackage[hidelinks]{hyperref}

\definecolor{bcnavy}{RGB}{11,31,58}
\definecolor{bcblue}{RGB}{32,74,135}
\definecolor{bcgold}{RGB}{176,141,87}
\definecolor{bcgrey}{RGB}{90,98,112}
\definecolor{bcrule}{RGB}{210,214,220}
\definecolor{bckeep}{RGB}{232,237,244}
\definecolor{bcact}{RGB}{247,241,228}
\definecolor{bcwarn}{RGB}{250,238,238}

\titleformat{\section}{\color{bcnavy}\Large\bfseries}{\thesection}{0.6em}{}
\titleformat{\subsection}{\color{bcblue}\large\bfseries}{\thesubsection}{0.6em}{}
\titlespacing*{\section}{0pt}{1.3em}{0.45em}

\pagestyle{fancy}\fancyhf{}
\renewcommand{\headrulewidth}{0pt}
\renewcommand{\footrulewidth}{0.4pt}
\renewcommand{\footrule}{\hbox to\headwidth{\color{bcrule}\leaders\hrule height \footrulewidth\hfill}}
\fancyfoot[L]{\footnotesize\color{bcgrey}Bayesian Capital}
\fancyfoot[C]{\footnotesize\color{bcgrey}Confidential --- Internal Research}
\fancyfoot[R]{\footnotesize\color{bcgrey}Page \thepage\ of \pageref{LastPage}}

\setlength{\parindent}{0pt}\setlength{\parskip}{0.5em}

\begin{document}

\begin{titlepage}
\vspace*{1.2cm}
{\color{bcgold}\rule{\textwidth}{2pt}}\\[0.4cm]
{\color{bcnavy}\Huge\bfseries Seven-Name Book}\\[0.25cm]
{\color{bcnavy}\Huge\bfseries Implementation Specification}\\[0.35cm]
{\color{bcblue}\LARGE Revision 3 --- Frozen Configuration}\\[0.5cm]
{\color{bcgold}\rule{\textwidth}{1pt}}\\[0.6cm]
{\large\color{bcgrey} Five names on the daily Bayesian / Ornstein--Uhlenbeck model}\\[0.2cm]
{\large\color{bcgrey} Two names on the weekly single-tranche model}\\[0.2cm]
{\large\color{bcgrey} All seven above the 50\% floor in the tested half; two fall short on the full sample}\\[1.2cm]
{\color{bcnavy}\large\bfseries Bayesian Capital}\\[0.15cm]
{\color{bcgrey} Quantitative Research}\\[0.15cm]
{\color{bcgrey} 4 August 2026 --- Revision 3, frozen}\\[1.6cm]

\fbox{\parbox{0.92\textwidth}{\color{bcnavy}\textbf{Purpose.}\\ \color{black}
The operative configuration for the reduced book. Four names --- TSLA, SPOT, SOFI and PLTR ---
are removed rather than repaired, having failed on both models. All seven retained names clear a
50\% annualised floor on the tested half of the sample; on the full sample MU and VRT fall short,
which \S7 addresses. NVDA and AVGO move from the daily model to a weekly single Monday
tranche. The remaining five keep their existing daily parameters unchanged.

\textbf{Revision 3 freezes this configuration.} The only change since Revision 2 is a 26-week
maximum holding period on the two weekly names, adopted as a safety rail: at that length it never
binds in this sample --- zero forced exits, returns identical to uncapped --- so it costs nothing
observed while bounding an otherwise open-ended exposure. Section 8 sets out why 26 rather than
the 12 weeks first indicated, and records the re-fit test run with the cap in place before
freezing. No parameter changes.

\textbf{Revision 2 restated every daily figure on the mark-to-market basis.} Revision 1 took each
name's terminal value from the model's Fund column, which holds its pre-purchase value while a
tranche is open, so a position still held at the end was counted at cost. Correcting it lowers
the daily names by 3 to 11 percentage points and the book's planning figure from 65\% to 59\%.
No parameter, weight or name selection changes as a result. MU's history has also been corrected
and extended to 3 August 2026. Section 6 sets out both changes.

All figures are on the execution-verified basis: every same-day round trip is checked against
minute or 5-minute bars, so no fill depends on an intraday ordering that cannot be known from
daily data.}}
\end{titlepage}

\section{The book}

\begin{table}[H]\centering\small
\begin{tabular}{llrrrr}
\toprule
\textbf{Name} & \textbf{Model} & \textbf{Verified} & \textbf{Tested half} & \textbf{Planning} &
\textbf{Weight}\\
\midrule
\rowcolor{bckeep} RKLB & Daily hybrid   & 101\% & 65\%  & 97\%  & 14.29\%\\
AVGO & Weekly Monday  & 91\%  & 96\%  & 60\%  & 14.29\%\\
\rowcolor{bckeep} NVDA & Weekly Monday  & 82\%  & 82\%  & 58\%  & 14.29\%\\
TSM  & Daily hybrid   & 63\%  & 89\%  & 51\%  & 14.29\%\\
\rowcolor{bckeep} VST  & Daily hybrid   & 61\%  & 55\%  & 54\%  & 14.29\%\\
VRT  & Daily hybrid   & 54\%  & 168\% & \textbf{44\%} & 14.29\%\\
\rowcolor{bckeep} MU   & Daily hybrid   & \textbf{43\%} & 124\% & 50\%  & 14.29\%\\
\midrule
\textbf{Book} & equal weight & \textbf{71\%} & --- & \textbf{59\%} & 100\%\\
\textbf{Book ex-RKLB} & & \textbf{66\%} & --- & \textbf{53\%} & \\
\bottomrule
\end{tabular}
\caption{Verified = full sample at the deployed configuration, marked to market and annualised
over each series' true span. Tested half = the portion scored outside every training window.
Planning = the figure to budget on, derived in \S5. Bold marks the two figures below the 50\%
floor.}
\end{table}

\begin{table}[H]\centering\small
\begin{tabular}{lrrr}
\toprule
\textbf{Name} & \textbf{Revision 1} & \textbf{Revision 2} & \textbf{Change}\\
\midrule
\rowcolor{bckeep} RKLB & 111\% & 101\% & $-10$pp\\
TSM  & 68\% & 63\% & $-5$pp\\
\rowcolor{bckeep} VST  & 65\% & 61\% & $-4$pp\\
VRT  & 65\% & 54\% & $-11$pp\\
\rowcolor{bckeep} MU   & 54\% & 43\% & $-11$pp\\
\midrule
NVDA, AVGO & \multicolumn{3}{l}{unchanged --- the weekly runner already valued open positions at
the close}\\
\bottomrule
\end{tabular}
\caption{Effect of the basis change on the daily names.}
\end{table}

\textbf{Two names no longer clear the 50\% floor on every measure.} MU returns 43\% on the full
sample and VRT plans at 44\%. Both clear it comfortably on the tested half --- 124\% and 168\% ---
which is the measure that decided the book's composition, but the unqualified claim that every
name clears 50\% no longer holds. \S7 sets out what to do about it.

\section{Daily names --- parameters}

Five names, unchanged from the current workbooks. Nine attempts to improve these parameters were
tested out of sample across this research and all were rejected; they are reproduced here exactly
as they stand. Cell references are to the \texttt{2 Tranche StopLoss 50d} sheet.

\begin{table}[H]\centering\footnotesize
\setlength{\tabcolsep}{4pt}
\begin{tabular}{lrrrrrr}
\toprule
& \multicolumn{6}{c}{\textbf{Bayes sleeve}}\\
\cmidrule(l){2-7}
\textbf{Name} & $\lambda$ (B2) & $\phi_L$ (D2) & $\psi$ (F2) & $k$ (H2) & prem (J2) &
peak cap (L2)\\
\midrule
\rowcolor{bckeep} RKLB & 0.352847 & 0.237406 & 0.113048 & 0.712293 & 0.026842 & 0.008015\\
TSM  & 0.462001 & 0.254287 & 0.065039 & 1.004530 & 0.015652 & 0.030484\\
\rowcolor{bckeep} VST  & 0.373520 & 0.264322 & 0.096187 & 1.385660 & 0.021816 & 0.013663\\
VRT  & 0.420699 & 0.425851 & 0.081656 & 1.116400 & 0.021691 & 0.048379\\
\rowcolor{bckeep} MU   & 0.600000 & 0.263600 & 0.010000 & 1.054800 & 0.024300 & 0.038000\\
\bottomrule
\end{tabular}
\end{table}

\begin{table}[H]\centering\footnotesize
\setlength{\tabcolsep}{5pt}
\begin{tabular}{lrrrrrr}
\toprule
& \multicolumn{4}{c}{\textbf{OU sleeve}} & \multicolumn{2}{c}{\textbf{Common}}\\
\cmidrule(lr){2-5}\cmidrule(l){6-7}
\textbf{Name} & $W$ (B3) & buffer $k$ (D3) & prem (H3) & cap (J3) & stop (T2) & comm (N2)\\
\midrule
\rowcolor{bckeep} RKLB & 85  & 0.220593 & 0.022147 & 0.033647 & 50 & 0.005\\
TSM  & 77  & 0.399229 & 0.010022 & 0.020893 & 50 & 0.005\\
\rowcolor{bckeep} VST  & 122 & 0.238599 & 0.014498 & 0.052746 & 50 & 0.005\\
VRT  & 80  & 0.567888 & 0.026603 & 0.059499 & 50 & 0.005\\
\rowcolor{bckeep} MU   & 91  & 0.300000 & 0.025000 & 0.029200 & 50 & 0.005\\
\bottomrule
\end{tabular}
\caption{Interest on idle cash (R2) is 3.14\% for all names. The stop period is a uniform 50
calendar days: a variable stop was tested, gained 8.9pp in sample and failed walk-forward at
13 folds of 30.}
\end{table}

\subsection{Sleeve split}

Bayes share \textbf{0.75} for all five names, entered flat rather than per name. Per-name Bayes
shares were tested and not supported: training windows selected a Bayes-heavy share for
essentially every name, so the per-name freedom bought nothing that the flat tilt did not
already provide.

\section{Weekly names --- parameters}

NVDA and AVGO run a \textbf{single Monday tranche} each. The three-tranche staggered version was
tested and the simplification costs nothing: NVDA scores 82.0\% on one tranche against 80.3\% on
three, and its parameter neighbourhood median \emph{improves} from 56.6\% to 58.6\%. AVGO gives
up 2.8pp of headline (90.7\% against 93.5\%) and gains on the neighbourhood, 64.8\% against
56.7\%. One tranche is both simpler and no worse.

\subsection{The rule}

Each Monday, if flat:
\begin{align*}
\text{buy} &= \min\bigl(\text{Monday open},\ \text{ATH}\times(1-\text{cap})\bigr)\\
\text{target} &= \text{buy} + \text{previous week close}\times\text{prem}
\end{align*}
where ATH is the running maximum of weekly highs through the previous week. Fill on the first
session of the week whose low reaches the buy price. Sell on any session at or after the fill
whose high reaches the target. If the target is not reached, \textbf{carry the position} and keep
the same target --- it is not recomputed while holding. Re-enter the Monday after a sale.
There is no stop loss.

\begin{table}[H]\centering\small
\begin{tabular}{lrrrrrr}
\toprule
\textbf{Name} & \textbf{cap} & \textbf{prem} & \textbf{Max hold} & \textbf{Anchor} &
\textbf{Tranches} & \textbf{Trades / yr}\\
\midrule
\rowcolor{bckeep} NVDA & 0.0935 & 0.0293 & 26 weeks & Monday & 1 & 20\\
AVGO & 0.0800 & 0.1000 & 26 weeks & Monday & 1 & 6.5\\
\bottomrule
\end{tabular}
\caption{Commission 0.005/share and 3.14\% interest on idle cash, as for the daily names. The
maximum hold never binds in this sample at 26 weeks; see \S8.}
\end{table}

\subsection{On the dormant mean-reversion term}

The original weekly workbook computes an entry level
\[
L=\min\Bigl(\operatorname{mean}\bigl(m\tfrac{H_{-1}+L_{-1}}{2},\ C_{-1}w\bigr)
+\log_{10}(H_{-1}-L_{-1})\,g,\ \text{Monday open}\Bigr)
\]
with $w=1.0082$, $m=1.1627$, $g=2.0374$. On the deployed values this term \textbf{never binds}:
it sits a median 9.4\% above the Monday open and was the binding constraint in 0 of 120 NVDA
weeks and 0 of 120 AVGO weeks. The entry is therefore always the Monday open or the ATH cap, as
written in \S3.1. Retain $w$, $m$ and $g$ if replicating the sheet, but they are inert ---
$g$ can be multiplied or divided by five with no change to a single trade. Do not lower $m$:
below roughly 1.0 the term switches on and performance falls (NVDA 80.3\% to 68.5\%, AVGO 53.4\%
to 37.3\%).

\section{Allocation sheet}

\begin{table}[H]\centering\small
\begin{tabular}{lll}
\toprule
\textbf{Setting} & \textbf{Cell} & \textbf{Value}\\
\midrule
\rowcolor{bckeep} Number of names & --- & \textbf{7}\\
Weight per name & --- & \textbf{0.142857} $(=1/7)$\\
\rowcolor{bcact} Floor: min weight per stock & B8 & \textbf{0.142857}\\
Cap: max weight per stock & E11:E17 & \textbf{0.142857}\\
\rowcolor{bckeep} Bayes fraction --- flat & B6 & \textbf{0.75}\\
Split mode & --- & \textbf{1} (flat, not per-name)\\
\rowcolor{bckeep} Interest on idle cash & R2 & \textbf{0.0314}\\
\bottomrule
\end{tabular}
\caption{Equal weighting is not a default --- it was tested against inverse-volatility,
Sharpe-proportional and return-proportional schemes and beat all of them by roughly 10pp.
Setting floor and cap both to $1/7$ pins the book to equal weight.}
\end{table}

The Bayes fraction and the weight apply only to the five daily names. The two weekly names take
their $1/7$ of capital into a single Monday tranche with no sleeve split.

\section{Performance --- what to budget}

\subsection{Why the planning figures sit below the verified ones}

Three discounts apply, each measured rather than assumed.

\textbf{The tilt is in-sample.} The Bayes share was walk-forwarded across 30 folds. It won 20 of
them and selected a Bayes-heavy share in all 30, so the direction is supported --- but the
realised out-of-sample gain was approximately \textbf{+2pp}, against the 7pp the in-sample curve
suggests (the five daily names average 57\% at 50/50 and 64\% at 75\%). Daily planning figures are
therefore the 50/50 result plus 2pp, not the 75\% result.

\textbf{The weekly parameters sit on a narrow ridge.} NVDA's 82\% is the peak of a region whose
median is 58.6\% and whose lower quartile is 45.4\%. AVGO's neighbourhood median is 64.8\% with a
lower quartile of 57.5\%. The neighbourhood median is what to budget, because a ridge fitted to
this sample is unlikely to stay exactly where this sample put it.

\textbf{The weekly figures below are uncapped.} At 26 weeks the cap makes no difference to
either name in this sample, so NVDA's 82.0\% and AVGO's 90.7\% stand unchanged, as do the
neighbourhood medians of 58.6\% and 64.8\% behind the planning figures.

\textbf{The weekly model ends fully invested.} Each weekly tranche finishes holding, so all of
its terminal value is an open position. It is already valued at the close here; marking it 20\%
lower would take NVDA to roughly 64\% and AVGO to 76\%.

\subsection{The half-sample test}

Weeks 1--59 sit inside every walk-forward training window; weeks 60--120 were scored outside at
least one. Splitting the two separates a configuration that works from one that fitted an early
run. This test is what removed four names, and the seven retained were held to the same standard.

\begin{table}[H]\centering\small
\begin{tabular}{lrrr}
\toprule
\textbf{Name} & \textbf{First half} & \textbf{Tested half} & \textbf{Shift}\\
\midrule
\rowcolor{bckeep} VRT  & $-12$\% & 168.0\% & $+180$pp\\
MU   & $-10$\% & 124.0\% & $+134$pp\\
\rowcolor{bckeep} TSM  & 40\%  & 89.0\%  & $+49$pp\\
AVGO & 72.2\%  & 96.4\%  & $+24$pp\\
\rowcolor{bckeep} NVDA & 44.2\%  & 81.8\%  & $+38$pp\\
VST  & 68\%  & 55.0\%  & $-13$pp\\
\rowcolor{bckeep} RKLB & 146\% & 65.0\%  & $-81$pp\\
\midrule
\rowcolor{bcwarn} \textit{Removed:} PLTR & 87.2\% & 6.6\% & $-81$pp\\
\rowcolor{bcwarn} \textit{Removed:} SPOT & 60.7\% & 3.6\% & $-57$pp\\
\rowcolor{bcwarn} \textit{Removed:} SOFI & 19.1\% & 13.0\% & $-6$pp\\
\rowcolor{bcwarn} \textit{Removed:} TSLA & $-0.2$\% & 19.2\% & $+19$pp\\
\bottomrule
\end{tabular}
\caption{All seven retained names clear 50\% in the tested half. VRT and MU are negative in the
untested half and strongest in the tested one, which is the opposite of a fitted result. The four
removed names are shown at the 50/50 split on which they were assessed.}
\end{table}

\subsection{Per-name detail}

\begin{table}[H]\centering\small
\begin{tabular}{lrrrrrr}
\toprule
\textbf{Name} & \textbf{50/50} & \textbf{75\% Bayes} & \textbf{Sharpe} & \textbf{max DD} &
\textbf{buys/yr} & \textbf{stops}\\
\midrule
\rowcolor{bckeep} RKLB & 95\% & 101\% & 1.16 & 46\% & 96 & 5\\
TSM  & 49\% & 63\% & 1.61 & 27\% & 65 & 3\\
\rowcolor{bckeep} VST  & 52\% & 61\% & 1.13 & 40\% & 68 & 5\\
VRT  & 42\% & 54\% & 0.95 & 62\% & 54 & 4\\
\rowcolor{bckeep} MU   & 48\% & 43\% & 0.93 & 42\% & 51 & 4\\
\bottomrule
\end{tabular}
\caption{Daily names, marked to market. MU is the one name where the tilt to 75\% Bayes reduces
return; its Sharpe also falls, so 50/50 would be defensible for MU alone. The book keeps a flat
share because per-name shares failed validation.}
\end{table}

\section{What changed since Revision 1}

\subsection{Mark-to-market accounting}

The model reports a terminal value from the Fund column, following the workbook's
\texttt{Y4 = Y867+AF867-P2}. That column keeps its pre-purchase value while a tranche is open, so
a position still held on the last day was counted at cost rather than at the closing price. Every
daily figure in Revision 1 carried that overstatement.

VST isolates the two components: it ends with nothing held, so its whole 4pp change is the switch
from the workbook's fixed 2.2-year annualisation to each series' true span. VRT and RKLB, which
end with both tranches holding, carry roughly a further 7pp of genuine open-position
overstatement on top of that.

\begin{table}[H]\centering\small
\begin{tabular}{lrrrr}
\toprule
\textbf{Name} & \textbf{Fund basis} & \textbf{Marked to market} & \textbf{Gap} &
\textbf{Tranches holding}\\
\midrule
\rowcolor{bckeep} VRT  & 65\%  & 54\%  & $+11$pp & 2 of 2\\
RKLB & 111\% & 101\% & $+11$pp & 2 of 2\\
\rowcolor{bckeep} MU   & 54\%  & 43\%  & $+11$pp & 1 of 2\\
TSM  & 68\%  & 63\%  & $+5$pp  & 1 of 2\\
\rowcolor{bckeep} VST  & 65\%  & 61\%  & $+4$pp  & 0 of 2\\
\bottomrule
\end{tabular}
\caption{MU's 11pp combines the basis change with its data correction, below.}
\end{table}

This is a reporting change, not a modelling one. No parameter, weight or name selection changes,
and the ordering of the names is unaffected. It does mean the workbook's own annual-return cells
read high whenever a tranche is open on the last row --- worth remembering when reading the sheet
rather than these tables.

\subsection{MU's price history}

MU's series ran to 24 June 2026 with a part-session value on that last day. It now runs to
3 August 2026, with 24 June rebuilt from the 5-minute bars (its low was 1031.00 against 991.10 in
the bars). The repair alone changed nothing --- 52\% and 54\% before and after --- but the added
27 sessions cover a fall from about 1233 to 739 and cost roughly 8pp. Over that month the model
lost 7.7\% while MU itself fell 20.9\%, so it cushioned the fall rather than avoiding it.

\section{What was tested and rejected}

The four removed names were each put through the full weekly pipeline before being dropped ---
own-parameter grid, walk-forward, neighbourhood and half-sample split.

\begin{table}[H]\centering\small
\begin{tabular}{lrrrl}
\toprule
\textbf{Name} & \textbf{Optimised} & \textbf{Neighbourhood} & \textbf{Walk-forward} &
\textbf{Reason for removal}\\
\midrule
\rowcolor{bcwarn} PLTR & 125.9\% & 95.9\% & 1/3 & 365\% first half, 3.1\% tested half\\
SOFI & 79.7\%  & 54.8\% & 1/3 & tested half only 23.3\%\\
\rowcolor{bcwarn} TSLA & 49.4\%  & 21.0\% & 2/3 & optimum on the grid boundary; folds disagreed\\
SPOT & 35.2\%  & 24.2\% & 0/3 & tested half $-29.1$\%\\
\bottomrule
\end{tabular}
\end{table}

Also tested and rejected on the retained names: five-parameter re-optimisation of the weekly model
(0 of 3 folds), the same with $g$ frozen (0 of 3), two-parameter re-optimisation (1 of 3), and
anchor mixing across weekdays (0 of 3, and it lowers the neighbourhood median from 56.6\% to
48.2\%). Across four experiments, re-fitting the weekly parameters has won 1 fold in 9.

The Monday anchor itself \emph{is} supported. Weekend repricing is 1.75\% against 1.41\% for a
normal overnight gap, and the effect tracks the number of non-trading days rather than the
weekday label --- a Tuesday following a three-day break reprices 1.88\% against 1.35\% for an
ordinary Tuesday. Monday is the best anchor for NVDA in every sub-period and for AVGO under both
its own and NVDA's parameters.

\section{The maximum holding period, and the freeze}

\subsection{Why a cap at all}

The weekly rule carries an unfilled target indefinitely. That is cheap in practice --- NVDA's
median hold is 4 days and AVGO's 20 --- but the exposure is formally unbounded, and it is the only
position in the book with no exit condition other than the target being reached. MU, tested for
the weekly model and rejected, reached 451 days in a single position. NVDA and AVGO top out at 118
and 122 days in this sample, so the tail is real but not extreme.

\subsection{Why 26 weeks and not 12}

Judged at the deployed parameters, a 12-week cap looked free: NVDA's tested half was unchanged at
81.8\%, AVGO's rose 1pp, and the cap fired only 8 times in 2.3 years across three names.

Judged on the \emph{parameter neighbourhood} --- the basis every planning figure in this
document rests on --- it is not free.

\begin{table}[H]\centering\small
\begin{tabular}{lrrrrr}
\toprule
& \multicolumn{2}{c}{\textbf{NVDA}} & \multicolumn{2}{c}{\textbf{AVGO}} &\\
\cmidrule(lr){2-3}\cmidrule(lr){4-5}
\textbf{Max hold} & at deployed & nbhd median & at deployed & nbhd median &
\textbf{Forced exits}\\
\midrule
\rowcolor{bcwarn} 8 weeks  & 59.4\% & 42.8\% & 76.7\% & 52.9\% & 3 / 4\\
\rowcolor{bcwarn} 12 weeks & 70.7\% & 46.3\% & 92.4\% & 63.9\% & 2 / 3\\
16 weeks & 80.0\% & 56.2\% & 84.6\% & 54.3\% & 1 / 2\\
\rowcolor{bckeep} 26 weeks & 82.0\% & 52.2\% & 90.7\% & 64.8\% & 0 / 0\\
none & 82.0\% & 58.6\% & 90.7\% & 64.8\% & 0 / 0\\
\bottomrule
\end{tabular}
\caption{A 12-week cap costs NVDA 12pp of neighbourhood median. At 26 weeks the cap does not bind
at the deployed parameters for either name.}
\end{table}

At 26 weeks nothing observed changes: zero forced exits, returns identical to uncapped, and
AVGO's neighbourhood median is unchanged at 64.8\%. NVDA's neighbourhood median still reads 6pp
below uncapped, because the cap binds at other points in the region even though it does not bind
at the deployed values; the series is also non-monotonic across lengths (56.2\% at 16 weeks,
54.8\% at 20, 52.2\% at 26), which is itself a sign that these differences are within the noise.

\textbf{The cap is therefore a safety rail, not a working rule.} It bounds an open-ended exposure
at six months without altering any historical behaviour. Do not tighten it without repeating this
test: shorter is not free.

\subsection{Re-fitting with the cap in place}

The cap was chosen with parameters frozen, which is the correct structural test but leaves the
parameters themselves fitted under a different rule. Before freezing, both names were re-fitted
with the cap active.

\begin{table}[H]\centering\small
\begin{tabular}{lrrrr}
\toprule
\textbf{Name} & \textbf{Deployed} & \textbf{Re-fitted} & \textbf{In sample} &
\textbf{Walk-forward}\\
\midrule
\rowcolor{bckeep} NVDA & 0.0935 / 0.0293 & 0.1400 / 0.0500 & $+12.7$pp & 1 of 3, $-19.9$pp\\
AVGO & 0.0800 / 0.1000 & 0.0800 / 0.1850 & $-0.2$pp & 0 of 3, $-26.7$pp\\
\bottomrule
\end{tabular}
\caption{NVDA's fitted peak cap scattered 5.6$\times$ across folds; AVGO's re-fit found nothing at
all in sample.}
\end{table}

One loose end was closed at the same time. Both earlier uncapped fits preferred a wider premium
for NVDA than the deployed 0.0293, and the time cap was the one change that could have justified
it, by bounding the wait that makes a wide target risky. Fitting the premium alone, with the peak
cap frozen, gives \textbf{0 of 3 folds and $-51.2$pp}, and the full-sample optimum lands at
0.0275 --- effectively the deployed value, which is also the best cell on the sweep at 70.7\%
against 52.4\% at 0.025 and 45.5\% at 0.035. The 0.055 the joint fit kept selecting only looked
good alongside a peak cap the folds could not agree on.

\subsection{Frozen}

Across six experiments, re-fitting the weekly parameters has now won \textbf{1 fold in 18}. The
configuration in \S2 and \S3 is frozen. The cap is the only change adopted since the parameters
were set, and it was adopted on a structural test with the parameters held fixed --- which this
re-fit confirms was the right way round.

\section{Risks and the 50\% floor}

\textbf{Two names miss the floor on the full sample.} MU returns 43\% and VRT plans at 44\%. Both
are emphatic on the tested half --- 124\% and 168\% --- and both were negative over the first
half, so the full-sample figure is an average across a period in which the configuration plainly
did not work and one in which it did. Three readings are available and the evidence does not
choose between them:
\begin{itemize}[leftmargin=1.4em,topsep=0.2em]
\item the tested half is the honest measure, both names pass, and nothing changes;
\item the full sample is the honest measure, and MU and VRT should be reviewed at the next
      quarter's data rather than dropped on 27 sessions of one name's history;
\item the floor was always a screening rule rather than an admission rule, and a book whose worst
      name plans at 44\% is not a worse book than one which reached 50\% by deletion.
\end{itemize}
The recommendation is the second: keep both, and revisit at the next refresh. Dropping MU on the
strength of one bad month, immediately after adding it, would be reacting to the most recent data
rather than to the weight of it.

\textbf{Concentration.} Four of the seven --- NVDA, AVGO, TSM and VRT --- are substantially the
same semiconductor and AI-infrastructure trade. The tested half was a strong period for that
complex, and the second-half strength of those four is partly one macro move counted four times.
This is the book's largest unhedged exposure and it increased when the four removed names left.
The diversifier shortlist was re-tested on the same half-sample basis and none of MRNA, OXY, FSLR,
COIN or DVN clears anything close to 50\% on the tested half, so the concentration cannot be fixed
without lowering the floor.

\textbf{Name-level selection.} The seven were identified as the survivors using the same 2024--26
sample on which they are scored. The half-sample test was run precisely to check this and they
passed it, but a 50\% floor met by removing everything below the line is a weaker claim than a
strategy that earns 50\% per name.

\textbf{RKLB dominance.} RKLB contributes roughly 14pp of the book's 59\% planning figure. Without
it the book plans at 53\%. Its tested half is 65\% against a first half of 146\%, the largest
decline of any retained name.

\textbf{No price stop on the weekly names.} A position is carried until the target is reached or
the 26-week cap is hit, whichever comes first. There is no loss-based exit. Median holding is 4
days for NVDA and 20 for AVGO; the longest in this sample are 118 and 122 days, both inside the
cap, so the cap has never yet been the binding exit.

\textbf{AVGO trade count.} A single Monday tranche produces roughly 6.5 completed trades a year on
AVGO. That is a thin base, and the name's optimum is a single grid cell that its neighbours do not
match --- 68\% at prem 0.08 and 64\% at 0.12 against 94\% at 0.10. AVGO is the least secure of the
seven.

\textbf{VRT's drawdown.} At 62\% peak-to-trough, VRT is the deepest of the five daily names by a
wide margin --- the next is RKLB at 46\%. Its Sharpe of 0.95 is the second lowest. VRT earns its
place on return, not on the path taken to it.

\vfill
{\color{bcrule}\rule{\textwidth}{0.6pt}}\\[0.2em]
{\footnotesize\color{bcgrey} All returns are execution-verified: same-day round trips are
confirmed against 1-minute bars for NVDA and 5-minute bars for the remaining names, with sessions
outside minute coverage resolved only where the fill is provable at the open. All are marked to
market, with any position open on the last day valued at that day's close, and annualised over
each series' true span. Commission 0.005 per share; idle cash earns 3.14\%.}

\end{document}
